% !TEX root = ../eigenfaces_algorithm.tex
%=====================================================
\section{PCA：找“差异最大”的那几个方向}

\subsection{目标函数：方差 = 信息}

我们想找单位向量 $\mathbf{u}$（$\norm{\mathbf{u}}=1$），
把所有残差脸投影到它上面后，投影值尽可能“散”。投影值是 $t_j=\inner{\mathbf{u}}{\mathbf{x}_j}=\mathbf{x}_j\mT\mathbf{u}$，
其平方和（中心化后均值自动为零，这就是方差的两倍尺度，差常数因子不影响结果）为
\[
\sum_{j=1}^{n} t_j^2
=\sum_j \mathbf{u}\mT\mathbf{x}_j\,\mathbf{x}_j\mT\mathbf{u}
=\mathbf{u}\mT\Big(\underbrace{\mat{X}\mT\mat{X}}_{d\times d}\Big)\mathbf{u}.
\]

\begin{intuition}
为什么“方差大 = 信息多”？想象一个班里所有人身高完全相同——身高这个维度方差为零，
告诉你身高也猜不出是谁。反过来，一个维度上方差很大，说明这个维度上人与人有明显差异，
\textbf{差异正是识别所需要的信息}。PCA 的哲学就是：把方差大的维度挑出来，
把方差小（近乎人人都一样）的维度扔掉。
\end{intuition}

\subsection{化成特征值问题（三行推导）}

带约束的最大化问题：
\[
\max_{\mathbf{u}}\ \mathbf{u}\mT\mat{X}\mT\mat{X}\,\mathbf{u}
\qquad\text{s.t.}\qquad \mathbf{u}\mT\mathbf{u}=1 .
\]
用 Lagrange 乘子法，令
\[
L(\mathbf{u},\lambda)=\mathbf{u}\mT\mat{X}\mT\mat{X}\mathbf{u}-\lambda(\mathbf{u}\mT\mathbf{u}-1),
\]
对 $\mathbf{u}$ 求导置零（利用 $\nabla_{\mathbf{u}}\,\mathbf{u}\mT\mat{C}\mathbf{u}=2\mat{C}\mathbf{u}$，$\mat{C}$ 对称）：
\[
2\,\mat{X}\mT\mat{X}\,\mathbf{u}-2\lambda\mathbf{u}=\mathbf{0}
\quad\Longrightarrow\quad
\boxed{\ \mat{X}\mT\mat{X}\,\mathbf{u}=\lambda\mathbf{u}\ }.
\]
结论漂亮得过分：
\begin{itemize}
  \item 让方差最大的方向，\textbf{恰好是} $\mat{X}\mT\mat{X}$ 的特征向量；
  \item 把它代回目标函数：$\mathbf{u}\mT(\lambda\mathbf{u})=\lambda$，
        \textbf{最大方差值恰好就是那个特征值}。
\end{itemize}
所以“按特征值从大到小排”，前 $k$ 个特征向量就是信息量最大的 $k$ 个方向。
这些 4096 维的特征向量 reshaped 回 $64\times 64$ 去看，就是传说中的
\textbf{特征脸（eigenfaces）}：前几张常像“整体明暗/光照方向”，后面几张像五官与纹理的细碎模式。

\begin{intuition}
把 $\mat{X}\mT\mat{X}$ 想成“像素与像素之间的相关性报告”。
它是对称阵，由谱定理知道它有一套互相垂直的拉伸方向。
PCA 做的事，说穿了就是\textbf{问这份报告：“你力气最大的拉伸方向是哪几个？”}
\end{intuition}
